% Generated draft: requires mathematical and source review.
\documentclass[11pt]{article}
\usepackage[margin=25mm]{geometry}
\usepackage{fontspec}
\setmainfont{Latin Modern Roman}
\usepackage{amsmath,amssymb,mathtools,unicode-math}
\setmathfont{Latin Modern Math}
\usepackage{enumitem,needspace,xurl,hyperref}
\setlength{\parindent}{0pt}
\setlength{\parskip}{5pt}
\setlength{\emergencystretch}{3em}
\newcommand{\N}{\mathbb{N}}
\newcommand{\Z}{\mathbb{Z}}
\newcommand{\Q}{\mathbb{Q}}
\newcommand{\R}{\mathbb{R}}
\newcommand{\C}{\mathbb{C}}
\newcommand{\F}{\mathbb{F}}
\newcommand{\A}{\mathbb{A}}
\newcommand{\PP}{\mathbb{P}}
\newcommand{\E}{\mathbb{E}}
\newcommand{\eps}{\varepsilon}
\newcommand{\dd}{\,\mathrm{d}}
\begin{document}
\section*{UnsolvedMath ID: 700026}
\textbf{Source problem code:} catalogue entry
\subsection{Definitions and mathematical statement}
Let \(G\) be an arbitrary group. Call \(G\) \emph{residually finite} if, for every \(g\in G\setminus\{1\}\), there exist a finite group \(F\) and a homomorphism \(\varphi:G\to F\) such that \(\varphi(g)\ne 1\). Interpret the \emph{maximum condition for subgroups} as the ascending chain condition: every chain \(H_1\le H_2\le\cdots\) of subgroups of \(G\) stabilizes, meaning that there exists \(n\ge 1\) such that \(H_m=H_n\) for every \(m\ge n\).

A group \(P\) is \emph{polycyclic} if it has a finite series
\[
\{1\}=P_0\trianglelefteq P_1\trianglelefteq\cdots\trianglelefteq P_r=P
\]
in which every quotient \(P_i/P_{i-1}\) is cyclic, either finite or infinite. Interpret \emph{almost polycyclic} as \emph{virtually polycyclic}: having a polycyclic subgroup of finite index. These are conventional interpretations of terminology that the supplied record does not explicitly define.

\textbf{Question (M. I. Kargapolov, 1965; Kourovka Notebook Problem 1.31).} Under these interpretations, is it true that, for every residually finite group \(G\) satisfying the maximum condition for subgroups, there exists a polycyclic subgroup \(P\le G\) such that
\[
[G:P]<\infty?
\]
\subsection{Short English statement}
Must every residually finite group whose ascending chains of subgroups stabilize contain a polycyclic subgroup of finite index? This uses the conventional meaning of ``almost polycyclic.''
\subsection{Sources}
\begin{itemize}
\item Local UnsolvedMath registry: supplied complete source record, permanent ID \(700026\), exact source problem code \texttt{catalogue entry}, titled \emph{Kourovka Notebook Problem 1.31}; upstream revision \texttt{d371ee917ac9647fbb06cd6fc529deec6f3dda82}.
\item The supplied attribution identifies \emph{The Kourovka Notebook (Unsolved Problems in Group Theory)}, original problem number 1.31, first issue (1965), as proposed by M. I. Kargapolov. The cited edition is the 21st issue edition (2026), \texttt{arXiv:1401.0300v48}. The supplied collection description names E. I. Khukhro and V. D. Mazurov as editors of Issues 1--21 (1965--2026).
\item The local record labels the problem \emph{open} and gives an update timestamp of \texttt{2026-10-16T14:00:00Z}. This is supplied registry metadata, not an independently verified assessment of the literature. The cited Notebook text has not been independently consulted for this draft.
\end{itemize}
\end{document}
